\documentclass[11pt]{article}

\usepackage[T1]{fontenc}
\usepackage[utf8]{inputenc}
\usepackage{amsmath,amssymb,amsthm}
\usepackage{booktabs}
\usepackage{geometry}
\usepackage{hyperref}
\usepackage{microtype}
\usepackage{xcolor}

\geometry{margin=1in}
\setlength{\emergencystretch}{2em}
\hypersetup{
  colorlinks=true,
  linkcolor=blue!55!black,
  citecolor=blue!55!black,
  urlcolor=blue!55!black
}

\newtheorem{theorem}{Theorem}[section]
\newtheorem{lemma}[theorem]{Lemma}
\newtheorem{proposition}[theorem]{Proposition}
\newtheorem{corollary}[theorem]{Corollary}
\theoremstyle{definition}
\newtheorem{definition}[theorem]{Definition}
\theoremstyle{remark}
\newtheorem{remark}[theorem]{Remark}

\newcommand{\F}{\mathbb F}
\newcommand{\bd}{\partial}
\newcommand{\Comp}{\operatorname{Comp}}
\newcommand{\dist}{\operatorname{dist}}
\newcommand{\HS}{\mathrm{HS}}
\newcommand{\symdiff}{\mathbin{\triangle}}
\numberwithin{equation}{section}

\title{A Radius-One Obstruction in Reconfiguring\\
Nowhere-Zero Flows toward Five-Cycle Double Covers}
\author{Atharva Vaidya\\
\small AI-assisted research draft for independent human verification}
\date{July 28, 2026}

\begin{document}
\maketitle

\begin{center}
\fcolorbox{red!65!black}{red!4}{
\parbox{0.91\textwidth}{
\textbf{Status and AI-use disclosure.}
The standard five-cycle double cover conjecture remains open.  This
manuscript proves a short packing-to-switch lemma and gives exact finite
counterexamples only to a proposed radius-one reconfiguration lemma.
It proves neither FiveCDC nor its negation.

OpenAI Codex agents using GPT-5-series models, under human direction,
formulated the auxiliary question and the packing-to-switch lemma, found
the finite examples, wrote the search and both checking programs,
performed the computations, drafted the proof and this manuscript, and
conducted an AI-assisted literature search.  The ``independent checker''
is independent of the discovery implementation, not independent of AI.
No human author, referee, or cited researcher is represented as having
verified these results.  Before submission, a human author must inspect
every proof, independently rerun and preferably reimplement the
enumerations, audit the graph and cover certificates, review the
literature, and accept ordinary scholarly responsibility.
}}
\end{center}

\begin{abstract}
Hu\v{s}ek and \v{S}\'amal recently gave a component-parity condition on
a nowhere-zero \(\F_2^3\)-flow which is equivalent, after choosing the
flow, to the existence of a five-cycle double cover.  Cranston, Li, Su,
Wang, and Xu study reconfiguration of nowhere-zero flows by changes
supported on cycles.  We test the direct radius-one synthesis of these
ideas.

First we prove a human-checkable bridge.  If one value-class matching of
a cubic \(\F_2^3\)-flow admits two edge-disjoint joins in its complement,
then an explicitly determined even support, and therefore a finite
sequence of legal simple-cycle switches, reaches a flow satisfying the
Hu\v{s}ek--\v{S}\'amal condition.  Connectedness of that support is
exactly what makes one switch sufficient.

We then give an exact 26-vertex obstruction to radius one.  The graph is
simple, cubic, cyclically \(4\)-edge-connected, has girth five, and is
not Tait-colourable.  One displayed value class already admits two
edge-disjoint boundary joins.  Nevertheless, for the displayed
nowhere-zero flow, a standard-library checker exhausts all \(9{,}213\)
simple cycles and all \(1{,}485\) legal cycle--value switches; none
reaches the component condition.  The packing construction splits into
two circuit components, and switching both reaches the condition.  Thus
the displayed flow has exact distance two.  An independently checked
five-cycle double cover of the same graph confirms that this is not a
counterexample to FiveCDC.

For calibration we record a cyclically \(4\)-edge-connected,
girth-five, Tait-colourable 40-vertex radius-one obstruction, and a
60-vertex positive control repaired by one \(25\)-edge cycle.  The
latter is also a literal instance of the packing-to-switch lemma.  All
finite claims are accompanied by solver-independent, certificate-based
checking code.  Possible novelty is stated conservatively and remains
subject to specialist prior-art review.
\end{abstract}

\section{Scope and relation to prior work}

A \emph{five-cycle double cover} (FiveCDC) of a finite graph is an
indexed list \(C_0,\ldots,C_4\) of Eulerian edge-subsets such that every
edge lies in exactly two members.  Eulerian means even degree at every
vertex; neither connectedness nor \(2\)-regularity is required.  Empty
members are allowed, so this is the usual ``at most five'' convention.
Only the standard, unoriented conjecture is considered here.

Hu\v{s}ek and \v{S}\'amal state FiveCDC as
Conjecture~1.9 and give an equivalent flow-selection formulation as
Conjecture~3.19 in \cite{HusekSamal2026}.  In particular, their
Theorem~3.16 supplies the component-parity criterion used below.  Thus
FiveCDC remains open, and the remaining universal issue is to select a
suitable nowhere-zero \(\F_2^3\)-flow.

Cranston, Li, Su, Wang, and Xu define adjacency of nowhere-zero flows by
requiring their difference to be supported on a cycle
\cite{CranstonEtAl2026}.  In the exponent-two group \(\F_2^3\), a
difference supported on a simple cycle has one constant nonzero value
around that cycle.  The switches below are precisely this special case,
with legality made explicit so that no zero edge is created.

Our contribution has deliberately narrow scope:
\begin{enumerate}
\item a direct proof that a two-join packing yields a finite legal switch
      sequence to a Hu\v{s}ek--\v{S}\'amal-good flow;
\item an exact strict 26-vertex countermodel to the assertion that one
      simple-cycle switch always suffices even when a value class already
      packs, together with an exact distance-two path and a positive
      FiveCDC certificate; and
\item one 40-vertex negative and one 60-vertex positive control.
\end{enumerate}
No reduction from FiveCDC to a universal radius-one assertion is claimed
here, and no statement about orientable cycle double covers is made.

\section{Flows, component defects, and switches}

Put \(W=\F_2^3\), written additively, and identify \(W\) with
\(\{0,\ldots,7\}\) by binary coordinates.  Let \(G\) be a finite simple
cubic graph.  A map
\[
                 f:E(G)\longrightarrow W-\{0\}
\]
is a nowhere-zero flow if
\[
                 \sum_{e\ni v}f(e)=0
                 \qquad(v\in V(G)).
\]
At a cubic vertex the three incident values are nonzero and pairwise
distinct.  Consequently each value class
\[
                 M_a(f)=\{e:f(e)=a\},\qquad a\ne0,
\]
is a matching.

For a nonzero linear functional \(\mu\in W^*\), identified with its
binary dot-product vector, define
\begin{equation}
\begin{split}
 F_\mu(f)&=\{e:\mu(f(e))=1\},\\
 K_\mu(f)&=\{e:\mu(f(e))=0\}=E(G)-F_\mu(f).
\end{split}
\end{equation}
For \(r\in W\) with \(\mu(r)=1\), let
\[
 Z_r(f)=\bd M_r(f)
\]
be the endpoints of the matching \(M_r(f)\).  Define
\begin{equation}
 d_\mu(f)=
 \#\left\{Q\in\Comp(V(G),K_\mu(f)):
                    |Q\cap Z_r(f)|\ \text{is odd}\right\}.
\end{equation}

\begin{lemma}[affine-value independence]\label{lem:affine}
The component parities in \emph{(2.2)} do not depend on the choice of
\(r\) satisfying \(\mu(r)=1\).
\end{lemma}

\begin{proof}
Fix a component \(Q\) of \(K_\mu(f)\).  The parity of
\(|Q\cap Z_r(f)|\) equals the parity of the number of \(r\)-valued edges
in \(\delta(Q)\), since an internal matching edge contributes two
endpoints.  Every edge in \(\delta(Q)\) has \(\mu\)-value one.

Write \(\ker\mu=\{0,p,q,p+q\}\), so the four affine values are
\(r,r+p,r+q,r+p+q\).  Let their boundary multiplicities modulo two be
\(n_0,n_p,n_q,n_{p+q}\).  Summing the flow equations on \(Q\) gives
\[
 n_0r+n_p(r+p)+n_q(r+q)+n_{p+q}(r+p+q)=0.
\]
Comparing coefficients in the basis \((r,p,q)\) gives
\[
 n_0+n_p+n_q+n_{p+q}=0,\quad
 n_p+n_{p+q}=0,\quad n_q+n_{p+q}=0.
\]
Hence all four multiplicities are equal.  This proves the claim.
\end{proof}

\begin{definition}
The flow \(f\) is \emph{H--S-good} if \(d_\mu(f)=0\) for some
\(\mu\ne0\).  Its defect profile is
\[
                       d(f)=(d_1(f),\ldots,d_7(f)).
\]
\end{definition}

By Theorem~3.16 and Observation~3.15 of
\cite{HusekSamal2026}, an H--S-good flow produces a standard FiveCDC.
Their Conjecture~3.19 asserts that every bridgeless cubic graph has some
H--S-good flow and is equivalent to FiveCDC.

\begin{definition}[legal simple-cycle switch]
Let \(a\in W-\{0\}\), and let \(C\) be a simple cycle of \(G\) such that
\(C\cap M_a(f)=\varnothing\).  Define
\begin{equation}
 f^{a,C}(e)=
 \begin{cases}
 f(e)+a,&e\in E(C),\\
 f(e),&e\notin E(C).
 \end{cases}
\end{equation}
\end{definition}

The cycle equation preserves flow conservation, and
\(C\cap M_a(f)=\varnothing\) is exactly the condition preventing a zero
value.  We write \(\dist_{\HS}(f)\) for the minimum number of legal
simple-cycle switches needed to reach an H--S-good flow, when such a
sequence exists.

\section{A human packing-to-switch lemma}

For an edge set \(J\), write \(\bd J\) for the set of vertices having odd
degree in \(J\).  A \(T\)-join is an edge set \(J\) with \(\bd J=T\).

\begin{theorem}[packing-to-switch]\label{thm:packing}
Fix \(a\ne0\) and \(\mu(a)=1\).  Put
\[
 F=F_\mu(f),\qquad M=M_a(f),\qquad J_0=F-M.
\]
Suppose \(G-M\) contains two edge-disjoint \(\bd M\)-joins
\(J_1,J_2\).  Then switching \(a\) on
\begin{equation}
                       C=J_0\symdiff J_1
\end{equation}
produces an H--S-good flow.  The support \(C\) is a disjoint union of
legal simple cycles, so this operation is a legal finite sequence of
simple-cycle switches.  If \(C\) is connected, one switch suffices.
\end{theorem}

\begin{proof}
The binary characteristic function \(\mu\circ f\) is a flow, so \(F\)
is Eulerian.  Because \(M\subseteq F\),
\[
                         \bd J_0=\bd(F-M)=\bd M.
\]
Thus \(J_0\) and \(J_1\) have the same boundary, and
\(C=J_0\symdiff J_1\) is Eulerian.  Both avoid \(M\), hence so does
\(C\).

In a cubic graph every vertex has degree zero or two in an Eulerian edge
set.  Every nonempty connected component of \(C\) is therefore a simple
cycle.  Switch \(a\) on all these components.  Since \(\mu(a)=1\), the
new affine support is
\begin{equation}
\begin{split}
 F\symdiff C
   &=F\symdiff\bigl((F-M)\symdiff J_1\bigr)\\
   &=M\cup J_1.
\end{split}
\end{equation}
The union is disjoint because \(J_1\subseteq E(G-M)\).

The matching \(M_a\) itself is unchanged: the support avoids \(M_a\), and
an edge with old value different from \(a\) cannot acquire value \(a\)
after adding \(a\), since that would force its old value to be zero.
Consequently every component switch remains legal when the components
are performed sequentially.

Finally, \(J_2\) avoids both \(M\) and \(J_1\), so by (3.2) it lies in
the new zero-side \(K_\mu\).  Restricting the equality
\(\bd J_2=\bd M\) to any component \(Q\) of this zero-side and applying
the handshaking lemma gives
\[
                       |Q\cap\bd M|\equiv0\pmod2.
\]
By Lemma~\ref{lem:affine}, this is the H--S condition for \(\mu\).
\end{proof}

\begin{remark}
The theorem exposes the radius-one gap.  A packing gives an Eulerian
switch support canonically, but that support may have several cycle
components.  The theorem guarantees a finite reconfiguration sequence;
it guarantees one step only when the support is connected.
\end{remark}

\section{The direct radius-one assertion}

The most direct strengthening suggested by
Theorem~\ref{thm:packing} is:

\begin{quote}
\emph{For every simple cubic cyclically \(4\)-edge-connected graph and
every nowhere-zero \(W\)-flow \(f\), either \(f\) is H--S-good or one
legal simple-cycle switch from \(f\) is H--S-good.}
\end{quote}

This statement concerns the current H--S component condition, not a
fixed five-cover and not an older all-value-class nonpacking condition.
The following exact countermodel refutes it.

\section{An exact strict 26-vertex countermodel}

Graph6 orders the edges lexicographically by second endpoint and then by
first endpoint.  Consider the graph
\begin{center}
\ttfamily\small
Y??G@cOOGC??P??Ap???`\_A??CIC???IC\_C?@??HO??A??B\_G????D?\_
\end{center}
and the following flow values in that edge order:
\begin{equation}
\begin{split}
(&5,1,5,4,4,4,7,1,6,7,3,1,6,4,3,2,3,2,6,1,\\
 &2,5,5,2,7,6,3,2,5,7,1,1,7,6,5,3,1,7,6).
\end{split}
\end{equation}

\begin{theorem}[strict radius-one countermodel]\label{thm:counter}
The graph and flow in \emph{(5.1)} have the following properties.
\begin{enumerate}
\item The graph is simple, cubic, cyclically \(4\)-edge-connected, has
      girth five, and has no Tait \(3\)-edge-colouring.
\item The displayed values form a nowhere-zero \(W\)-flow with profile
      \[
                           (6,6,4,6,2,4,4).
      \]
\item Its value-4 class packs two edge-disjoint boundary joins in its
      complement.
\item The graph has \(9{,}213\) simple cycles and exactly \(1{,}485\)
      legal cycle--value switches from this flow.  None is H--S-good.
      The minimum defect at radius at most one is two.
\item The graph has the explicit standard FiveCDC given below.
\item The displayed flow has exact H--S distance two.
\end{enumerate}
\end{theorem}

\paragraph{FiveCDC certificate.}
Encode a pair of cover coordinates by a five-bit integer.  In graph6
edge order, the labels are
\begin{equation}
\begin{split}
(&3,20,5,17,3,5,6,3,3,5,10,6,3,5,3,6,5,9,10,12,\\
 &6,18,24,12,20,6,18,6,20,18,6,6,18,20,6,5,5,6,3).
\end{split}
\end{equation}
Every label in (5.2) has Hamming weight two.  For each of the five bit
positions and each vertex, an even number of incident labels contain
that bit.  Thus the five bit-coordinate edge sets are Eulerian and every
edge belongs to exactly two of them.

\paragraph{Exact distance-two certificate.}
Starting from (5.1), perform
\begin{equation}
\begin{array}{c@{\qquad}l}
 a=4&
 C=(0,1,2,8,9,20,22,24,25,26,30,31,33),\\[2mm]
 a=4&
 C=(6,12,14,15,34,35).
\end{array}
\end{equation}
Both supports are simple cycles and avoid the current value class being
added.  The successive profiles are
\begin{equation}
                (4,6,4,2,4,4,4),\qquad(4,6,4,0,6,4,4).
\end{equation}
Thus (5.3) gives an H--S-good flow after two steps.

\paragraph{Packing certificate.}
The value-4 matching is
\[
                         M_4=(3,4,5,13).
\]
The two edge sets
\begin{align*}
J_0={}&(1,14,15,18,20,21,26,28,29,30,31,32,35,37,38),\\
J_1={}&(0,2,6,7,8,9,12,16,17)
\end{align*}
are disjoint, avoid \(M_4\), and both have boundary
\(\partial M_4\).  The even support furnished by
Theorem~\ref{thm:packing} has exactly the two components in (5.3).
Hence this certificate explains the distance-two repair while proving
that packability alone does not force a radius-one neighbour.

\begin{proof}[Certificate proof of Theorem~\ref{thm:counter}]
The supplied standard-library checker decodes graph6 from first
principles.  It verifies distinct edges, degree three, connectedness,
the 26 vertex flow equations, and all 39 nonzero values.  It examines every deletion
of one, two, or three edges and rejects a deletion as a cyclic cut only
after constructing all components and confirming that at least two
contain cycles.  It independently computes girth and exhaustively
backtracks over proper three-edge-colourings.

For completeness of the radius-one enumeration, the checker performs a
depth-first enumeration of simple cycles.  The minimum vertex of a cycle
is its unique allowed starting vertex.  The two possible traversal
directions are quotiented by retaining the direction whose second vertex
is smaller than its final vertex.  Hence each undirected simple cycle is
generated exactly once.  For each generated cycle and each
\(a\in W-\{0\}\), the checker tests legality directly and, when legal,
recomputes the full seven-entry defect profile from graph components.
The resulting exact totals are \(9{,}213\) cycles, \(1{,}485\) legal
switches, and zero clean neighbours.  This proves
\(\dist_{\HS}(f)\geq2\).  Direct verification of (5.3)--(5.4) proves
\(\dist_{\HS}(f)\leq2\).

Finally, the checker verifies the matching and both join boundaries,
and tests every weight and every vertex-coordinate parity in (5.2).
These finite checks prove all six assertions.
\end{proof}

\begin{remark}
Theorem~\ref{thm:counter} is a counterexample only to the radius-one
auxiliary statement.  Its explicit FiveCDC proves that it is not a
FiveCDC counterexample.
\end{remark}

\section{Forty- and sixty-vertex controls}

Two controls separate radius-one failure from the existence of a cover
and from the packing-to-switch mechanism.

\subsection{A Tait-colourable 40-vertex negative control}

There is a simple cubic, cyclically \(4\)-edge-connected graph of girth
five with a nowhere-zero flow having profile
\begin{equation}
                         (6,8,10,4,10,6,6).
\end{equation}
The checker exhausts \(941{,}438\) simple cycles and \(48{,}544\) legal
cycle--value switches.  None is H--S-good, and the minimum defect at
radius at most one is two.  In contrast to
Theorem~\ref{thm:counter}, this graph has a displayed Tait
three-edge-colouring; its three colour-pair \(2\)-factors give a
three-cycle double cover.  This is therefore a clean negative control
for the reconfiguration assertion, not for FiveCDC.

\subsection{A 60-vertex positive control and the packing lemma}

The 60-vertex control is simple cubic and cyclically
\(4\)-edge-connected.  Its starting profile is
\begin{equation}
                         (8,8,4,6,6,8,6).
\end{equation}
Switching \(a=7\) on the \(25\)-edge simple cycle
\begin{equation}
\begin{split}
C=(1,2,8,11,18,22,24,26,27,30,31,33,35,37,39,40,\\
\hspace{31mm}42,43,45,62,65,67,73,74,75)
\end{split}
\end{equation}
gives profile
\begin{equation}
                         (10,6,6,8,8,6,0).
\end{equation}
Thus this state is repaired in one step.

The repair also instantiates Theorem~\ref{thm:packing}.  For the
value-\(7\) matching, two edge-disjoint \(\bd M_7\)-joins in
\(G-M_7\) are
\begin{equation}
\begin{split}
J_1={}&(1,2,13,16,18,20,26,32,35,44,46,49,50,53,57,60,61,\\
      &\qquad62,63,64,69,71,72,73,75,81,83),\\
J_2={}&(3,6,7,10,12,14,17,19,22,24,27,29,34,37,40,41,43,\\
      &\qquad45,47,48,55,56,74,80,82,85,86,87,88).
\end{split}
\end{equation}
With \(F=F_7(f)\), the checker verifies
\begin{equation}
             C=(F-M_7)\symdiff J_1,\qquad
             J_2\cap(M_7\cup J_1)=\varnothing.
\end{equation}
Equations (6.3)--(6.6) are therefore a literal connected-support case
of the human lemma.

\begin{table}[t]
\centering
\caption{The three principal exact computations.  ``Legal'' counts
cycle--value pairs, not merely cycles.}
\begin{tabular}{@{}rrrrrl@{}}
\toprule
order & cycles & legal & initial minimum & radius-one minimum & outcome\\
\midrule
26 & \(9{,}213\)   & \(1{,}485\)  & \(2\) & \(2\) & packs; no one-switch repair\\
40 & \(941{,}438\) & \(48{,}544\) & \(4\) & \(2\) & no repair; Tait control\\
60 & ---           & ---          & \(4\) & \(0\) & displayed one-switch repair\\
\bottomrule
\end{tabular}
\end{table}

\section{Verification architecture}

The principal checker is
\[
\texttt{scratch/verify\_husek\_samal\_packable\_one\_switch\_countermodel.py}.
\]
It uses only the Python standard library.  Its graph6 decoder, flow
checker, component-defect routine, simple-cycle generator, switch
routine, cyclic-cut enumeration, girth computation, Tait-colouring
backtracker, and FiveCDC checker are implemented locally.  No SAT solver
or opaque mathematical library is used.

The 40-vertex negative control is checked by the separately written
\[
\texttt{scratch/verify\_husek\_samal\_one\_switch\_boundary.py}.
\]

The 60-vertex construction has a second checker,
\[
\texttt{scratch/verify\_fano\_order60\_flow\_repair.py}.
\]
It reconstructs the graph from two 32-vertex pole recipes, verifies the
one-switch repair and the two joins, and separately follows a
13-switch path to the flow induced by an explicit FiveCDC.  The two
programs share no imported project module.

The frozen artifact hashes at the time of this draft are:
\begin{center}
\scriptsize
\begin{tabular}{@{}p{0.70\textwidth}p{0.25\textwidth}@{}}
\toprule
SHA-256 & file\\
\midrule
\nolinkurl{9bdcb3a65f6dc07224f8a20c79e45e9085c9f73c55ead745ac64127decd76607}
& packable-state checker\\
\nolinkurl{ac51aaf266178eb764fcada85b25e17dc2c18c129413abb51a2b309513ea415d}
& boundary checker\\
\nolinkurl{09f9bdc797948b49cf6b7a3080160e07111dad718f0e07ae00a2a9c0958b9a6e}
& independent 60-vertex replay\\
\nolinkurl{5cd280e730e0dc336681c3768945b7c691e023cb4f957a7408ea306597053f45}
& packable order-26 retained state\\
\nolinkurl{6e1dc5fb3b290fe8364559844ed419ee6dda1ec974b2ad3aebfd21442c8fe6aa}
& order-40 retained state\\
\bottomrule
\end{tabular}
\end{center}

From the project root, the exact replay commands are
\begin{verbatim}
python3 scratch/verify_husek_samal_packable_one_switch_countermodel.py
python3 scratch/verify_husek_samal_one_switch_boundary.py
python3 scratch/verify_fano_order60_flow_repair.py
\end{verbatim}

\section{What remains open}

The finite countermodels show only that radius one is not a universal
domination bound for H--S-good flows, even in a strict reduced-looking
class.  They leave open all of the following:
\begin{enumerate}
\item whether every nowhere-zero \(\F_2^3\)-flow can reach an H--S-good
      flow by some finite legal cycle-switch sequence;
\item whether a larger universal radius exists;
\item whether every bridgeless graph has some initially H--S-good flow,
      the equivalent formulation of FiveCDC in
      \cite{HusekSamal2026}; and
\item the standard FiveCDC conjecture itself.
\end{enumerate}

The 26-vertex state has exact distance two, but this is one instance,
not evidence of a universal radius-two theorem.  It also shows sharply
that the finite sequence in the packing-to-switch lemma cannot always be
replaced by one circuit.  The lemma supplies no universal packing
theorem.

\section{Novelty and responsibility boundary}

The cited papers establish the H--S component criterion and the general
cycle-supported reconfiguration framework.  We make no novelty claim
for either.  The elementary packing-to-switch bridge, the exact
radius-one formulation, and the three finite certificates may be new,
but that assessment is provisional.  The search and literature review
were AI-assisted and not exhaustive in the scholarly sense.  Priority
and publishability require review by specialists in cycle double covers,
nowhere-zero flows, and flow reconfiguration.

The human part of Theorem~\ref{thm:packing} is fully displayed.  The
finite theorems reduce to explicit data and short, inspectable exhaustive
programs, but those programs and this prose were generated with AI
assistance.  Reproducibility is not peer review, and implementation
independence is not human independence.

\bibliographystyle{plain}
\bibliography{references}

\end{document}
